\section{Methods}\label{sec:methods}
% length (W3, trim pass + second fixer pass): about 1,165 whitespace tokens of prose on non-comment lines against the 1,000-word budget (first draft 1,912, fixer pass 1,586, trim pass 1,082). Episode detection thresholds, augmentation probabilities, the reference learning rate, bag sizes, fold pool hours and the MIL training details are kept in comments next to the prose they belong to; restore any of them if space allows. Details that Section V states itself (seeds, the gated-attention ablation, permutation importance on the validation split) are not repeated here. The fold pool sizes stay in the comment under Part C.

All methods read only the 20 official channels of Section~\ref{sec:dataset} and quantities derived from them. Parts A--B and F are the first round, Parts C--E and G the pre-registered campaign.

\subsection{Window Models}

A window is a $600 \times 20$ array (60 s at 10 Hz), standardised per channel with training-split statistics. We compare four sequence classifiers \cite{wang2017tsc,ismailfawaz2019review}: a 1-D ResNet \cite{wang2017tsc} on the 10 Hz signal and three token models on a shared two-layer convolutional stem with stride 4 (150 tokens of width 128): two bidirectional GRU layers \cite{cho2014gru}, a four-layer Transformer encoder \cite{vaswani2017attention} and four bidirectional Mamba-2 layers \cite{gu2023mamba,dao2024mamba2}, a pure-PyTorch implementation of the published algorithm. All four share one pooling head whose hidden vector is the window embedding reused in Part E. Parameters: ResNet 2.31 M, bi-GRU 0.73 M, Transformer 0.70 M, Mamba-2 1.07 M.
% inputs clipped, missing values set to 0; Mamba-2 block = the chunked state-space-duality scan run forward and on the reversed tokens, without custom kernels

\subsection{Reference Training Recipe}

Training uses class-balanced sampling weights $(N / (K n_k))^{0.5}$ capped at 10 ($N$ windows, $K$ classes, $n_k$ per class), AdamW \cite{loshchilov2019adamw} with warm-up over the first 3\%
of the steps and cosine decay, batches of 256, label smoothing 0.05 and early stopping on validation window macro-F1. The feature baselines feed 180 window statistics (nine per channel) and the physical features proposed by large language models (LLMs) in Part F to a class-balanced logistic regression and a histogram gradient-boosting classifier \cite{pedregosa2011sklearn,ke2017lightgbm}.
% at most 25 epochs, patience 6; learning rate 1e-3 (5e-4 for the Transformer) and weight decay 0.05 (experiments/train.py DEFAULTS and LR_BY_MODEL; not in OUTLINE.md); statistics = moments, extremes, quantiles, roughness and coverage; features and statistics are evaluated alone and together, under the same splits and protocols; seeds are stated in Section V

\subsection{Model-Disjoint Internal Folds and Pre-Registration}

Because the released validation split shares car models with training (Section~\ref{sec:dataset}), every campaign choice was made on internal car-model-disjoint folds drawn per task from the train + validation pool: three folds per task hold out 19 of 57 candidate car models (size) and 16, 15 and 20 (BEV vs ICE). A fold's training set drops every drive whose identity group also drove a held-out model; the one shared device pool spanning 19 car models is not an owner and loses only its drives of held-out models (fleet rule).
% pools (OUTLINE.md): size 168,588 windows, 2,810 h, 117 car models; BEV vs ICE 121,381 windows, 2,023 h, 107 car models; fleet group = 33 vehicles, about 480 h, 13.3 % of the size pool's windows

For the reference recipe, early stopping on the shared-model validation split selects epochs 16--22, on the folds 1--8 (median about 4): that validation rewards memorising car models. Final models therefore train for the median fold best epoch (GRU S9 11, ResNet S9 16, GRU B5 2, ResNet B5 7, episode model 10).
% the validation split is only reported (Section V)

Six gates were pre-registered; G0 and G5 are procedural (stored numbers reproduced; one frozen test read per final model). G1: the reference recipe's fold drive-level score lies within 5 points of its test value, else the folds would have been redesigned. G2: the best recipe reaches a fold floor of 70 (size) or, for BEV vs ICE, 90 and at least 1 point over the reference recipe (or 92 on human-only drives); on failure, size reports ceilings instead of a 90 target and BEV vs ICE keeps the reference recipe. G3: the episode model gains 2 points over the best window recipe (size) or reaches 91 (BEV vs ICE), else it is only an ablation. G4: an ensemble or stacker needs 1.0 points over the best single model out of fold (outcomes in Section~\ref{sec:experiments} and Table~\ref{tab:folds}).

\subsection{Recipe Changes S1--S9 and B1--B5}

Each ladder adds one change at a time to the reference S0 / B0 (S1--S6 cumulative, S7--S9 each on S5; the BEV-vs-ICE ladder skips S3 and S6). S1 / B1: car-model-balanced sampling inside each class (shares $\propto$ hours$^{0.5}$). S2 / B2: regression head on the log footprint. S3: CORAL ordinal head \cite{cao2020coral}. S4 / B3: log curb-mass and four-class powertrain heads (B3 also trains the auxiliary heads on HEV / PHEV windows). S5 / B4: physics-safe augmentation (joint sign flip of the six lateral channels, one masked 3--15 s span with the lead validity flag cleared, dropout of one channel group, Gaussian noise 0.02). S6: same-class cutmix across car models. S7 / B5: down-weighting of windows with driving automation system (DAS) activity by $(1 - \beta) + \beta (1 - s)$, $s$ the DAS share, $\beta = 0.5$. S8: conditional adversarial car-model head \cite{ganin2016dann,long2018cdan} with gradient reversal ($\lambda = 0.1$). S9: group-DRO-lite \cite{sagawa2020dro} (per-car-model loss re-weighting). Decisions stay softmax arg-max; the final recipes are S9 (size) and B5 (BEV vs ICE).
% augmentation details (OUTLINE.md): sign flip p 0.5, masked span p 0.5, channel groups = lead / steering torque / pedals with p 0.1; no speed or amplitude scaling and no time reversal (both change the vehicle signature); heads: SmoothL1 on the standardised log footprint; loss weights classification 1, CORAL 0.5, footprint 1, mass 0.5, powertrain-4 0.3; S9 weights = softmax over running per-car-model losses

\subsection{Episode-Centric Multiple-Instance Model}

The episode MIL (multiple-instance learning) model treats each drive as a bag of the transients in which mass and powertrain should show, trained for both tasks jointly. Eight episode types are detected on Savitzky-Golay-smoothed speed \cite{savitzky1964}: ACC (acceleration), DEC\_LIFT and DEC\_BRK (deceleration with the brake mostly released or pressed), STOP, LAUNCH, LEAD\_ACC and LEAD\_DEC (lead acceleration or deceleration onsets in car-following range) and TURN (sustained low-speed curvature); 735,571 episodes come from 8,662 drives (545 have none). An episode is an 18-s onset-anchored crop of 22 channels (the 20 inputs, smoothed acceleration, validity) with a 9-dimensional context vector; acceleration and deceleration episodes are also re-indexed on a speed grid (0--36 m/s, 0.5 m/s bins) as an acceleration-against-speed profile $a(v)$.
% detection thresholds left out for length (OUTLINE.md, Episodes and the MIL model): filter 5 samples / order 2; ACC / DEC = monotone segments (tolerance 0.3 m/s) of >= 3 s with |dv| >= 2 m/s and mean |a| >= 0.1 m/s^2, 2 s of cruise trimmed at both ends; DEC_LIFT = brake pressed in <= 10 % of the samples; STOP = standstill >= 1 s reached from >= 2 m/s within 15 s; LAUNCH = departure from standstill; LEAD_ACC / LEAD_DEC = lead slope >= 0.5 / <= -0.6 m/s^2 held 1 s after 2 s below 0.3 m/s^2 in car-following range (headway <= 6 s), ego response a >= 0.3 m/s^2 for 0.5 s within 8 s; TURN = |curvature| >= 0.01 1/m for >= 1 s at 2-12 m/s
% representation details left out for length (OUTLINE.md): context = anchor speed, speed change, duration, mean and peak acceleration, automation share, lead share, headway, gap flag; 180 stored samples, onset at sample 40 (STOP: standstill onset at sample 140); grid = 72 bins x 6 channels (smoothed acceleration, pedal flags, automation flag, validity, time from onset)

A shared encoder ($d = 128$; convolutional stem, residual blocks, type embedding, context MLP) embeds the crops, a grid encoder the $a(v)$ grid, and an adapter the frozen 60-s window embeddings of the window models (Part A) as a ninth instance type; on the folds the adapter uses only the fold's own window checkpoints. Gated attention \cite{ilse2018mil} pools each bag per task; heads predict size, BEV vs ICE, four-class powertrain, log footprint, log curb mass and two ordinal thresholds. The loss is the bag cross-entropies + 0.3 instance cross-entropies + 0.5 Huber(log footprint) + 0.25 Huber(log mass) + 0.25 cross-entropy (powertrain-4) + 0.25 binary cross-entropy (ordinal); about 0.99 M parameters train for a fixed 10 epochs (AdamW, warm-up, cosine decay). The bag prediction is the drive-level output; window-mean aggregation of the same model gives the row comparable with the window recipes.
% training details (OUTLINE.md unless marked): bag cross-entropies with label smoothing 0.05, masked by task eligibility; bags of up to 64 episodes (lead and turn episodes first, each capped at 16) + 8 window embeddings per drive in training, 128 + 16 at evaluation; batches of 24 drives with a cap on the share of one car model (code: 10 %, not in OUTLINE.md); augmentation = lateral sign flip p 0.5, instance dropout 0.2, lead dropout 0.1, noise 0.02, crop jitter (code: 160 of the 180 stored samples, offset 10 at evaluation, jitter 0-20 samples; not in OUTLINE.md); learning rate 1e-3 and weight decay 0.05 (experiments/mil_train.py DEFAULTS); the gated-attention vs learned-mean ablation is reported in Section V

\subsection{LLM-Assisted Feature Discovery}

Following Section~\ref{sec:related}, LLMs only propose programs that a verifier scores \cite{romeraparedes2024funsearch}; they never classify. Six proposer agents, one physics lens each, saw the channel definitions and a training-only sample of 2,727 windows and proposed 74 candidate features; six verifier agents ran the code, measured class separation and dropped 21, and an auditor removed 19 more, leaving 34 registered features. The coverage floor (the minimum share of windows on which a feature is defined) was lowered from 40\% to 10\%
because gradient boosting treats a missing regime as information; everything was fixed before any evaluation split was touched.
% pipeline details (OUTLINE.md): lenses = longitudinal powertrain, transmission signature, braking / regeneration / coasting, lateral inertia and geometry, car-following interaction, speed-dependent capability; 62 features + 12 from a gap-filling round; verifiers dropped 21 (coverage below the floor, degenerate or duplicate); the auditor's 19 = 6 in the first audit + 13 in the second (duplicates, a v^2 speed confound, controller-driven cues such as lead-response gain and lag and braking-onset time-to-collision, recording artefacts); the first pass with the 40 % floor let 12 features reach the auditor, the 10 % floor re-admitted 29; permutation importance is measured on the validation split (Section V)

\subsection{Aggregation to Drive and Vehicle Level}

Drive and vehicle protocols are those of Section~\ref{sec:dataset} (mean of the per-window probability vectors, arg-max). Out of fold, weighted (DAS, episode, steering) arithmetic and log-odds means, robust and gated means and a two-stage vehicle rule were compared with the plain mean on the reference recipe; none gained the required 1.0 drive-level points on both tasks (Section~\ref{sec:experiments}). A supplementary rule fixed before the final test run, the log-odds mean with window weights $(1 - s + 0.1)^{0.5} \times (1 + n_{\mathrm{acc}} + n_{\mathrm{dec}})$ ($s$ DAS share, $n$ acceleration / deceleration episode counts), is reported beside every final model.
